% Wersja źródłowa: audyt II, jednostka 31
\setcounter{section}{30}
\section{Algebry Banacha, widmo i holomorficzny rachunek funkcyjny}
\label{sec:banach-algebras31}

\providecommand{\atlasref}[2]{#2}

\noindent\textbf{Wyniki wykorzystywane.}
Korzystamy z twierdzeń Hahna--Banacha, Banacha--Alaoglu,
lematu Neumanna (Lemat~\ref{lem:neumann15}) oraz podstawowych
twierdzeń skalarnej analizy zespolonej (Cauchy'ego, Liouville'a,
o usuwaniu osobliwości i zasadzie maksimum). Twierdzenie
Stone'a--Weierstrassa udowodnimy w tym rozdziale.
\par\medskip

W poprzednim rozdziale wprowadziliśmy widmo ograniczonego operatora
na zespolonej przestrzeni Banacha, przede wszystkim w związku z operatorami
zwartymi. Teraz rozpatrzymy widmo elementu \emph{dowolnej} zespolonej algebry
Banacha. Otrzymamy w ten sposób wspólne podstawy teorii operatorów, algebr
funkcyjnych oraz algebr $C^*$. Nie powtarzamy twierdzenia Riesza--Schaudera:
jego twierdzenia będą przykładami ogólnego rachunku spektralnego.

\subsection{Algebra Banacha i elementy odwracalne}

\noindent\emph{Powiązania tematyczne w Bibliotece matematycznej:} \atlasref{LA-03::prop:lin3-end-algebra}{LA-03: algebra endomorfizmów}.


\begin{definicja}[Zespolona algebra Banacha]
\label{def:banach-algebra31}
\emph{Zespoloną algebrą Banacha} nazywamy zespoloną przestrzeń Banacha
$A$ z łącznym, dwuliniowym mnożeniem $A\times A\to A$ spełniającym
$\|ab\|\le\|a\|\|b\|$ dla wszystkich $a,b\in A$.
W tym rozdziale, o ile nie zaznaczono inaczej, zakładamy, że istnieje
element jednostkowy $1_A$ i $\|1_A\|=1$. Algebra nie musi być przemienna.
Element $a\in A$ nazywamy odwracalnym, gdy istnieje $b\in A$ takie,
że $ab=ba=1_A$; oznaczamy wtedy $b=a^{-1}$.
\end{definicja}

\begin{przyklad}[Trzy podstawowe algebry]
\label{ex:banach-algebras31}
Jeśli $X$ jest niezerową zespoloną przestrzenią Banacha, to $\mathcal L(X,X)$
z normą operatorową i składaniem jest algebrą Banacha.
Jeśli $K$ jest niepustą zwartą przestrzenią Hausdorffa, to $C(K,\C)$
z normą supremum i mnożeniem punktowym jest przemienną algebrą Banacha.
Przestrzeń $\ell^1(\N_0)$ z iloczynem splotowym
$(a*b)_n=\sum_{j=0}^n a_jb_{n-j}$ jest przemienną algebrą Banacha,
której jednostką jest $(1,0,0,\ldots)$.
W ostatnim przypadku
\[
\sum_{n\ge0}|(a*b)_n|\le
\sum_{n\ge0}\sum_{j=0}^n|a_j||b_{n-j}|
=\sum_{j,k\ge0}|a_j||b_k|=\|a\|_1\|b\|_1.
\]
Zmiana kolejności w sumach skończonych dla każdej
współrzędnej daje łączność splotu, a zamiana $j$ na
$n-j$ daje przemienność. Zupełność wynika z zupełności
przestrzeni $\ell^1$.
\end{przyklad}

\begin{lemat}[Szereg Neumanna i otwartość grupy odwracalnych]
\label{lem:neumann31}
Niech $A$ będzie zespoloną algebrą Banacha z jedynką
$1_A$ taką, że $\|1_A\|=1$.
Jeżeli $u\in A$ i $\|u\|<1$, to $1_A-u$ jest odwracalny i
$(1_A-u)^{-1}=\sum_{n=0}^\infty u^n$ (zbieżność w normie).
Zbiór $A^\times$ wszystkich elementów odwracalnych jest otwarty,
a odwzorowanie $a\mapsto a^{-1}$ jest na nim ciągłe.
\end{lemat}
\begin{proof}
Wersja operatorowa szeregu Neumanna została już udowodniona
w Lemacie~\ref{lem:neumann15}. Sprowadzimy do niej przypadek
dowolnej algebry Banacha z jedynką. Dla $v\in A$ określmy
operator lewego mnożenia $L_v:A\to A$, $L_v(x)=vx$.
Z podmultiplikatywności normy wynika $\|L_v\|\le\|v\|$,
a z łączności mnożenia $L_vL_w=L_{vw}$ i $L_{1_A}=I_A$,
gdzie $I_A$ jest operatorem identycznościowym na przestrzeni
Banacha $A$. Ponadto $L_v(1_A)=v$, więc odwzorowanie
$v\mapsto L_v$ jest iniektywne.

Jeżeli $\|u\|<1$, to $\|L_u\|<1$. Z Lematu~\ref{lem:neumann15}
otrzymujemy $B=(I_A-L_u)^{-1}=\sum_{n=0}^\infty L_u^n$
w normie operatorowej. Ewaluacja operatora na $1_A$ jest
ciągła, ponieważ $\|T(1_A)\|\le\|T\|\|1_A\|$.
Stąd $v=B(1_A)=\sum_{n=0}^\infty u^n$ w normie $A$.
Ciągłość $w\mapsto L_w$ i równość $L_{u^n}=L_u^n$
dają $L_v=B$. Wobec tego
$L_{(1_A-u)v}=I_A=L_{v(1_A-u)}$, a iniektywność
reprezentacji daje $(1_A-u)v=v(1_A-u)=1_A$.
Jest więc $v=(1_A-u)^{-1}$, jak należało wykazać.

Dla $a\in A^\times$ i $b\in A$ spełniającego
$\|a^{-1}(b-a)\|<1$ rozkład $b=a(1_A+a^{-1}(b-a))$
oraz udowodniona właśnie wersja szeregu dają odwracalność
$b$. To dowodzi otwartości $A^\times$, ponieważ wystarczy
$\|b-a\|<1/\|a^{-1}\|$. Jeśli
$\|a^{-1}(b-a)\|\le1/2$, oszacowanie szeregu daje
$\|b^{-1}\|\le2\|a^{-1}\|$ (przy $\|1_A\|=1$).
Tożsamość $b^{-1}-a^{-1}=b^{-1}(a-b)a^{-1}$ prowadzi
więc do $\|b^{-1}-a^{-1}\|\le2\|a^{-1}\|^2\|b-a\|$,
co dowodzi ciągłości odwracania.
\end{proof}

\subsection{Widmo i rezolwenta}

\begin{definicja}[Widmo elementu algebry]
\label{def:spectrum31}
Dla $a\in A$ definiujemy \emph{zbiór rezolwenty}
$\rho_A(a)=\{z\in\C:z1_A-a\in A^\times\}$ oraz
\emph{widmo} $\sigma_A(a)=\C\setminus\rho_A(a)$.
Na $\rho_A(a)$ definiujemy \emph{rezolwentę}
$R_a(z)=(z1_A-a)^{-1}$.
Indeks dolny $A$ pomijamy, gdy algebra jest ustalona. W przypadku
$A=\mathcal L(X,X)$ jest to definicja z poprzedniego rozdziału.
\end{definicja}

\begin{twierdzenie}[Zwartość i niepustość widma]
\label{thm:spectrum-nonempty31}
Niech $A$ będzie zespoloną algebrą Banacha z jedynką i $a\in A$.
Wtedy $\sigma_A(a)$ jest niepustym zwartym podzbiorem
$\{z\in\C:|z|\le\|a\|\}$. Rezolwenta $R_a$ jest holomorficzna
na $\rho_A(a)$ jako funkcja o wartościach w przestrzeni Banacha $A$.
Dla $z,w\in\rho_A(a)$ zachodzi tożsamość rezolwenty
\[
 R_a(z)-R_a(w)=(w-z)R_a(z)R_a(w).
\]
\end{twierdzenie}
\begin{proof}
Dla $|z|>\|a\|$ lemat~\ref{lem:neumann31} daje
$R_a(z)=\sum_{n\ge0}a^n z^{-n-1}$. Zatem widmo jest ograniczone,
a jego dopełnienie otwarte jako przeciwobraz $A^\times$
przez ciągłe odwzorowanie $z\mapsto z1_A-a$.
Dla $z_0\in\rho(a)$ mamy faktoryzację
$z1_A-a=(z_0 1_A-a)(1_A+(z-z_0)R_a(z_0))$.
Stąd dla $|z-z_0|\|R_a(z_0)\|<1$ otrzymujemy
\[
R_a(z)=\sum_{n\ge0}(-1)^n(z-z_0)^n R_a(z_0)^{n+1}.
\]
Szereg jest jednostajnie zbieżny w normie na każdym
mniejszym domkniętym dysku i dowodzi holomorficzności.
Dla tożsamości rezolwenty obliczamy bez założenia przemienności
algebry:
\[
R_a(z)[(w1_A-a)-(z1_A-a)]R_a(w)
=R_a(z)-R_a(w).
\]
Lewa strona jest równa $(w-z)R_a(z)R_a(w)$.

Gdyby widmo było puste, $R_a$ byłaby całkowitą funkcją holomorficzną
$\C\to A$. Dla $|z|>2\|a\|$ mielibyśmy
$\|R_a(z)\|\le 2/|z|$. Na pozostałym domkniętym dysku rezolwenta jest ograniczona
z ciągłości (możemy użyć dysku o promieniu $1+2\|a\|$).
Jest więc ograniczona na całej płaszczyźnie.
Dla dowolnego $\varphi\in A'$ funkcja skalarna
$\varphi\circ R_a$ jest całkowita i ograniczona, a z twierdzenia
Liouville'a stała. Granica przy $|z|\to\infty$ daje
$\varphi(R_a(z))=0$. Funkcjonały ciągłe rozdzielają punkty $A$,
więc $R_a(z)=0$, co przeczy $(z1_A-a)R_a(z)=1_A$.
\end{proof}

\begin{przyklad}[Widmo operatora mnożenia]
\label{ex:multiplication-spectrum31}
Dla $f\in C(K,\C)$ zachodzi $\sigma_{C(K)}(f)=f(K)$.
Istotnie, $z1-f$ ma odwrotność w $C(K)$ dokładnie wtedy,
gdy nigdzie nie znika; odwrotnością jest wówczas $1/(z-f)$.
W szczególności $\sigma(f)$ może być dowolnym niepustym zwartym
podzbiorem $\C$ (weźmy $K$ równy temu zbiorowi i $f(z)=z$).
\end{przyklad}

\begin{uwaga}[Widmo zależy od algebry]
Jeżeli $B\subseteq A$ jest domkniętą podalgebrą z tą samą jednostką,
to $\sigma_A(b)\subseteq\sigma_B(b)$ dla $b\in B$:
odwrotność w $B$ jest również odwrotnością w $A$.
Równość nie musi zachodzić. Niech
$\mathbb D=\{z\in\C:|z|<1\}$ i
$\mathbb T=\{z\in\C:|z|=1\}$.
Rozważmy algebrę dyskową
$\mathcal A(\overline{\mathbb D})$, tj. funkcje ciągłe
na $\overline{\mathbb D}$ i holomorficzne w $\mathbb D$.
Ograniczanie do brzegu daje izometryczną podalgebrę
$B\subseteq C(\mathbb T)$: izometrię zapewnia zasada maksimum,
a domkniętość wynika z zupełności algebry dyskowej:
ciąg Cauchy'ego w normie supremum na brzegu jest, na mocy
zasady maksimum, ciągiem Cauchy'ego na całym domkniętym
dysku. Jego granica jest ciągła i holomorficzna wewnątrz
dysku dzięki zbieżności jednostajnej na zbiorach zwartych.
Dla funkcji współrzędnej $u(z)=z$ na $\mathbb T$ zachodzi
$\sigma_{C(\mathbb T)}(u)=\mathbb T$.
Jednak $\sigma_B(u)=\overline{\mathbb D}$: dla
$|\lambda|>1$ funkcja $1/(u-\lambda)$ jest w $B$,
a dla $|\lambda|<1$ nie mogłaby być odwrotnością
$u-\lambda$ w $B$, gdyż jej holomorficzne przedłużenie
musiałoby spełniać $(z-\lambda)g(z)=1$ w dysku:
funkcja $(z-\lambda)g(z)-1$ znika na brzegu, więc
na mocy zasady maksimum znika w całym dysku.
Podstawienie $z=\lambda$ daje sprzeczność.
Dla $|\lambda|=1$ odwrotność nie istnieje już w
$C(\mathbb T)$.
\end{uwaga}

\begin{twierdzenie}[Wzór na promień spektralny]
\label{thm:spectral-radius31}
Niech $A$ będzie zespoloną algebrą Banacha z jedynką, $a\in A$.
Wówczas granica $\lim_{n\to\infty}\|a^n\|^{1/n}$ istnieje
i jest równa promieniowi spektralnemu
\[
 r(a):=\max\{|z|:z\in\sigma_A(a)\}
 =\lim_{n\to\infty}\|a^n\|^{1/n}
 =\inf_{n\ge1}\|a^n\|^{1/n}.
\]
\end{twierdzenie}
\begin{proof}
Nierówność $\|a^{n+m}\|\le\|a^n\|\|a^m\|$ daje istnienie granicy
$L$ i równość z infimum. Rzeczywiście, jeśli $a^m=0$ dla
pewnego $m$, teza o granicy jest natychmiastowa.
W przeciwnym razie, ustalając $m\ge1$ i zapisując $n=qm+r$,
$0\le r<m$, otrzymujemy
$\|a^n\|\le\|a^m\|^q\max_{0\le j<m}\|a^j\|$.
Po wyciągnięciu pierwiastka stopnia $n$ i przejściu do granicy
z góry dostajemy
$\limsup_n\|a^n\|^{1/n}\le\|a^m\|^{1/m}$.
Ponieważ $m$ jest dowolne, a każdy wyraz ciągu pierwiastków
jest nie mniejszy niż ich infimum, otrzymujemy żądaną granicę.
Jeżeli $|z|>L$, wybieramy $q$ z $L/|z|<q<1$.
Dla dostatecznie dużych $n$ mamy $\|a^n z^{-n}\|\le q^n$.
Szereg $\sum_{n\ge0}a^n z^{-n-1}$ jest więc normowo
zbieżny. Pomnożenie jego sumy częściowej z dowolnej strony
przez $z1_A-a$ daje $1_A-a^{N+1}z^{-N-1}$,
którego granicą jest $1_A$. Dowodzi to odwracalności.
Zatem $r(a)\le L$.

Dla $R>\|a\|$ z szeregu Neumanna i skalarnego wzoru Cauchy'ego
(stosowanego po złożeniu z dowolnym $\varphi\in A'$) otrzymujemy
\[
 a^n=\frac{1}{2\pi i}\int_{|z|=R}z^nR_a(z)\,\mathrm dz.
\]
Całka jest granicą sum Riemanna w zupełnej przestrzeni Banacha $A$;
krzywa ma dodatnią orientację. Holomorficzność rezolwenty pozwala
przesunąć okrąg do dowolnego promienia $R>r(a)$, gdyż pierścień
między okręgami nie przecina widma. W konsekwencji
$\|a^n\|\le R^{n+1}\sup_{|z|=R}\|R_a(z)\|$.
Po wyciągnięciu pierwiastka stopnia $n$ uzyskujemy $L\le R$.
Przechodząc z $R$ do $r(a)$ dostajemy tezę.
\end{proof}

\begin{przyklad}[Element quasi-nilpotentny]
Operator Volterry $V$ na $C([0,1],\C)$ z
poprzedniego rozdziału ma $\sigma(V)=\{0\}$,
mimo że $V\ne0$. Zatem $r(V)=0$ i
$\lim_n\|V^n\|^{1/n}=0$. Jednak dla funkcji stałej $1$ mamy
$(V^m1)(t)=t^m/m!$ z indukcyjnego wzoru na całkę
(Przykład~\ref{ex:volterra-spectrum30}), więc $V^m\ne0$
dla każdego $m\ge1$. Operator nie jest nilpotentny.
\end{przyklad}

\subsection{Charaktery, ideały maksymalne i transformacja Gelfanda}

\noindent\emph{Powiązania tematyczne w Bibliotece matematycznej:} \atlasref{AA-09::aa:def:prime-max-ideal}{AA-09: ideały maksymalne}.


\begin{definicja}[Charakter i przestrzeń charakterów]
\label{def:character31}
Niech $A$ będzie \emph{przemienną} zespoloną algebrą Banacha
z jedynką. Charakterem nazywamy niezerowy homomorfizm
zespolonych algebr $\chi:A\to\C$ (a priori bez założenia ciągłości).
Zbiór wszystkich charakterów oznaczamy $\Delta(A)$.
\emph{Transformacja Gelfanda} przyporządkowuje elementowi $a\in A$
funkcję $\widehat a:\Delta(A)\to\C$, $\widehat a(\chi)=\chi(a)$.
\end{definicja}

\begin{lemat}[Twierdzenie Gelfanda--Mazura]
\label{lem:gelfand-mazur31}
Każda zespolona algebra Banacha z jedynką, w której każdy niezerowy
element jest odwracalny, jest izomorficzna jako unitalna algebra z $\C$.
\end{lemat}
\begin{proof}
Dla $a\in A$ widmo $\sigma(a)$ jest niepuste na mocy
twierdzenia~\ref{thm:spectrum-nonempty31}. Wybierzmy
$z\in\sigma(a)$. Element $a-z1_A$ nie jest odwracalny,
więc przy założeniu o $A$ musi być równy zeru.
Zatem każdy $a$ jest skalarną wielokrotnością $1_A$.
\end{proof}

\begin{twierdzenie}[Charaktery i widmo elementu]
\label{thm:gelfand-spectrum31}
Niech $A$ będzie przemienną zespoloną algebrą Banacha z jedynką.
Wówczas każdy charakter $\chi\in\Delta(A)$ jest ciągły,
$\|\chi\|=1$, a $\Delta(A)$ jest niepustą zwartą przestrzenią
w topologii słabej-* $\sigma(A',A)$. Dla każdego $a\in A$ zachodzi
\[
 \sigma_A(a)=\{\chi(a):\chi\in\Delta(A)\},\qquad
 \|\widehat a\|_\infty=r(a)\le\|a\|.
\]
W szczególności $a\mapsto\widehat a$ jest ciągłym homomorfizmem
unitalnym $A\to C(\Delta(A),\C)$.
\end{twierdzenie}
\begin{proof}
Dla charakteru $\chi$ wybieramy $b\in A$ z $\chi(b)\ne0$.
Równość $\chi(b)=\chi(1_A)\chi(b)$ daje $\chi(1_A)=1$.
Jeżeli $\chi(a)\notin\sigma(a)$, element $a-\chi(a)1_A$
byłby odwracalny, chociaż jego obraz przez $\chi$ jest zerem.
To niemożliwe, więc $|\chi(a)|\le r(a)\le\|a\|$.
Zatem $\chi$ jest ciągły i ma normę $1$.

Jeżeli $z\in\sigma(a)$, element $a-z1_A$ jest nieodwracalny.
Ideał $A(a-z1_A)$ jest właściwy, bo w przemiennej
algebrze zawieranie jednostki oznaczałoby odwracalność
$a-z1_A$. Lemat Zorna stosujemy do właściwych ideałów
zawierających ten ideał: suma łańcucha jest ideałem
właściwym, ponieważ nie zawiera jednostki. Otrzymujemy
ideał maksymalny $M$. Jego domknięcie jest ideałem
z ciągłości mnożenia. Gdyby $\overline M=A$, istniałby
$m\in M$ z $\|m-1_A\|<1$, a więc odwracalny, co przeczy
właściwości $M$. Maksymalność daje $\overline M=M$.
Iloraz $A/M$ jest zupełny w normie ilorazowej; norma
jest podmultiplikatywna, bo można niezależnie przybliżać
infima norm reprezentantów obu klas. Jednostka $[1_A]$
ma normę nie większą niż $1$, a z niezerowości i
$[1_A]^2=[1_A]$ wynika nierówność odwrotna.
Dla $b\notin M$ maksymalność daje $M+Ab=A$, czyli
$1_A=m+cb$ dla pewnych $m\in M,c\in A$.
W ilorazie $[c][b]=[1_A]$; przemienność daje też
odwrotność prawostronną. Zatem $A/M$ jest zespoloną
algebrą Banacha, której każdy niezerowy element jest odwracalny. Z lematu~\ref{lem:gelfand-mazur31}
$A/M\simeq\C$, a odwzorowanie ilorazowe złożone z tym
izomorfizmem daje charakter $\chi$ spełniający $\chi(a)=z$.
To dowodzi formuły na widmo i niepustości $\Delta(A)$.

Topologia słaba-* czyni $\chi\mapsto\chi(a)$ ciągłym.
Zbiór charakterów jest w kuli jednostkowej $A'$ domknięty,
ponieważ równania $\chi(1_A)=1$ i
$\chi(ab)=\chi(a)\chi(b)$ dla wszystkich $a,b\in A$
określają przecięcie zbiorów domkniętych w topologii
słabej-*. Każdy funkcjonał w tym przecięciu jest
niezerowym homomorfizmem. Z twierdzenia Banacha--Alaoglu
kula jest zwarta, a topologia słaba-* jest Hausdorffa,
gdyż różne funkcjonały różnią się na pewnym elemencie $A$.
Norma $\widehat a$ jest maksimum $|\chi(a)|$, czyli $r(a)$.
Liniowość, multiplikatywność i zachowanie jednostki
transformacji wynikają z tych samych własności charakterów;
oszacowanie normy daje jej ciągłość.
\end{proof}

\begin{uwaga}[Algebry niepółproste]
Transformacja Gelfanda nie musi być różnowartościowa.
Na przykład w dwuwymiarowej algebrze
$A=\C[\epsilon]/(\epsilon^2)$, wyposażonej w normę
$\|a+b\epsilon\|=|a|+|b|$, istnieje tylko jeden charakter,
$\chi(a+b\epsilon)=a$. Niezerowy element $\epsilon$ przechodzi na
funkcję zerową. Jądro transformacji Gelfanda jest przecięciem jąder
wszystkich charakterów. Każde takie jądro jest ideałem
maksymalnym, gdyż charakter jest surjekcją na $\C$.
Odwrotnie każdy ideał maksymalny jest jądrem charakteru,
jak wykazano w Twierdzeniu~\ref{thm:gelfand-spectrum31}.
Przecięcie wszystkich ideałów maksymalnych to radykał
Jacobsona algebry przemiennej. Ze wzoru
$\|\widehat a\|_\infty=r(a)$ wynika, że składa się on
dokładnie z elementów o promieniu spektralnym zero.
\end{uwaga}

\subsection{Holomorficzny rachunek funkcyjny}

\noindent\emph{Powiązania tematyczne w Bibliotece matematycznej:} \atlasref{LA-18::thm:lin18-analytic-calculus}{LA-18: analityczny rachunek funkcyjny macierzy}.


\begin{definicja}[Rachunek funkcyjny Dunforda]
\label{def:holomorphic-calculus31}
Niech $A$ będzie zespoloną algebrą Banacha z jedynką, $a\in A$,
a $U\subseteq\C$ otwartym otoczeniem $\sigma_A(a)$.
Dla $f\in\mathcal O(U)$ wybieramy skończoną, zorientowaną
rodzinę zamkniętych krzywych kawałkami $C^1$, oznaczoną $\Gamma$,
leżących w $U\setminus\sigma(a)$ i stanowiących dodatnio
zorientowany brzeg obszaru wielokątnego $D$, dla którego
$\sigma(a)\subset D\subseteq\overline D\subset U$.
Dopuszczamy wiele składowych i dziury; wspólne krawędzie
kwadratów w konstrukcji brzegu znoszą się z przeciwnymi
orientacjami. Indeks $\operatorname{Ind}_\Gamma(w)$
wynosi $1$ w $D$ poza brzegiem i $0$ poza $\overline D$.
Takie kontury istnieją: wybieramy $d>0$ tak, aby zamknięte
$d$-otoczenie zwartego widma mieściło się w $U$.
Na siatce kwadratowej o średnicy oczka mniejszej niż
$d/4$ bierzemy wszystkie domknięte kwadraty przecinające
zamknięte $d/2$-otoczenie widma. Ich suma jest skończona,
zawiera widmo we wnętrzu i mieści się w $U$.
Wewnętrzne krawędzie usuwamy, a pozostałe orientujemy
tak, by wnętrze leżało po lewej stronie. W każdym
wierzchołku liczba wchodzących krawędzi jest równa
liczbie wychodzących (tak było dla każdego kwadratu
i pozostaje po skreśleniu par). Skończony zbiór tych
krawędzi rozkłada się więc na zamknięte drogi wielokątne,
również gdy niektóre kwadraty stykają się tylko wierzchołkami.
Dla punktu poza wszystkimi krawędziami siatki indeks
sumy brzegów wynika z addytywności indeksu: wynosi $1$
w jednym z wybranych kwadratów i $0$ poza nimi.
Dla punktu na usuniętej krawędzi lub wewnętrznym wierzchołku
ten sam wniosek wynika z lokalnej stałości indeksu poza
pozostawionym konturem.
Tę konstrukcję można powtórzyć w dowolnym mniejszym
otoczeniu widma, otrzymując kontury zagnieżdżone.
Definiujemy
\[
 f(a):=\frac{1}{2\pi i}\int_\Gamma f(z)(z1_A-a)^{-1}\,\mathrm dz.
\]
Całka oznacza granicę sum Riemanna po parametryzacjach konturów
w normie przestrzeni Banacha $A$; zapis obejmuje sumę całek
po wszystkich składowych konturu. Ciągłość funkcji podcałkowej na każdym odcinku parametryzacji
zapewnia istnienie całki w zupełnej przestrzeni $A$.
\end{definicja}

\begin{twierdzenie}[Własności holomorficznego rachunku funkcyjnego]
\label{thm:holomorphic-calculus31}
Niech $A,a,U$ będą jak w definicji~\ref{def:holomorphic-calculus31}.
Odwzorowanie $f\mapsto f(a)$ z $\mathcal O(U)$ do $A$ jest unitalnym
homomorfizmem algebr. Dla wielomianu $p$ otrzymujemy zwykły element
$p(a)$, a dla $z\notin\sigma(a)$ rachunek funkcyjny funkcji
$w\mapsto(z-w)^{-1}$ daje $R_a(z)$. Wartość $f(a)$ zależy wyłącznie
od kiełka $f$ w otoczeniu widma i jest zgodna z ograniczaniem $U$.
Ponadto
\[
 \sigma_A(f(a))=f(\sigma_A(a)).
\]
\end{twierdzenie}
\begin{proof}
Dla dwóch dopuszczalnych obszarów $D_1,D_2$ wybieramy
trzeci $D_0$ z $\sigma(a)\subset D_0$ i
$\overline D_0\subset D_1\cap D_2$. Funkcja
$z\mapsto\varphi(f(z)R_a(z))$ jest holomorficzna
na otoczeniu każdego $\overline D_j\setminus D_0$
dla $j=1,2$. Skalarne twierdzenie Cauchy'ego na tych
obszarach (z uwzględnieniem orientacji wewnętrznego
brzegu) daje równość całek po $\partial D_j$ i
$\partial D_0$. Ciągły funkcjonał komutuje z całką,
co wynika z sum Riemanna. Twierdzenie Hahna--Banacha
przenosi te równości na równość w $A$.
Ten sam argument dla wspólnego mniejszego otoczenia
dowodzi zgodności z ograniczaniem $U$ i zależności od kiełka.
Liniowość jest bezpośrednia. Dla funkcji całkowitych
$1$ i $z$ wolno najpierw rozszerzyć dziedzinę do $\C$
(zgodność z ograniczaniem), a następnie użyć brzegu dużego kwadratu zawierającego widmo,
na którym $|z|>\|a\|$, i rozwinięcia rezolwenty. Jednostajna zbieżność szeregu na tym brzegu pozwala
całkować wyraz po wyrazie. Całka z $z^k$ podzielona
przez $2\pi i$ wynosi $1$ dla $k=-1$ oraz $0$
dla pozostałych całkowitych $k$ (istnieje funkcja
pierwotna $z^{k+1}/(k+1)$ dla $k\ne-1$).
Dostajemy odpowiednio $1_A$ i $a$.

Dla mnożenia wybieramy obszary $D_i,D_o$ z
$\sigma(a)\subset D_i$, $\overline D_i\subset D_o$
i $\overline D_o\subset U$. Kontury $\Gamma_i,\Gamma_o$
są rozłączne; $z$ oznacza zmienną na $\Gamma_o$,
a $w$ na $\Gamma_i$. Ciągłość mnożenia i jednostajna
zbieżność sum Riemanna uzasadniają zapis iloczynu całek
jako całki podwójnej oraz zmianę kolejności całkowania.
Używamy tożsamości
$R_a(z)R_a(w)=(R_a(w)-R_a(z))/(z-w)$.
Po podstawieniu do podwójnej całki dla $f(a)g(a)$
i całkowaniu najpierw po konturze zewnętrznym w wyrazie z $R_a(w)$
otrzymujemy przez wzór Cauchy'ego $f(w)g(w)R_a(w)$;
wyraz z $R_a(z)$ znika, ponieważ odpowiadająca mu funkcja
$w\mapsto g(w)/(z-w)$ jest holomorficzna wewnątrz konturu
wewnętrznego. Istotnie, $\operatorname{Ind}_{\Gamma_o}(w)=1$
dla $w\in\Gamma_i$, a $z\in\Gamma_o$ nie należy do
$\overline D_i$. Skalarne wzory Cauchy'ego stosujemy
na wszystkich składowych z ich orientacjami, również
gdy $U$ jest niespójny lub obszary mają dziury.
Daje to $(fg)(a)=f(a)g(a)$, a wraz z liniowością
i obliczeniem $1(a),z(a)$ także zgodność z wielomianami.

Dla funkcji $g(w)=(z-w)^{-1}$, $z\notin\sigma(a)$,
rozważanej na dostatecznie małym otoczeniu widma, zachodzi $(z1_A-a)g(a)=1_A$ na mocy własności homomorfizmu;
Analogicznie $g(a)(z1_A-a)=1_A$ (obrazy funkcji
komutują). Stąd $g(a)=R_a(z)$.
Jeśli $\lambda\notin f(\sigma(a))$, funkcja
$1/(f-\lambda)$ jest holomorficzna na pewnym otoczeniu widma.
Z własności homomorfizmu wynika odwracalność
$f(a)-\lambda1_A$, czyli $\sigma(f(a))\subseteq f(\sigma(a))$.
Jeśli $\lambda=f(\mu)$ dla $\mu\in\sigma(a)$,
funkcja $h(w)=(f(w)-f(\mu))/(w-\mu)$ ma usuwalną
osobliwość w $\mu$ i jest holomorficzna na otoczeniu widma.
Gdyby $f(a)-f(\mu)1_A=(a-\mu1_A)h(a)$ było odwracalne,
to czynniki, które komutują, miałyby odwrotności
(jeżeli $bc=cb$ i $d=(bc)^{-1}$, to $b,c$ komutują
z $bc$, a więc także z $d$; stąd $b(cd)=(cd)b=1_A$
oraz $c(bd)=(bd)c=1_A$). Otrzymalibyśmy sprzeczność
z $\mu\in\sigma(a)$; stąd przeciwne zawieranie.
\end{proof}

\begin{przyklad}[Eksponenta i funkcje wymierne]
Jeśli $a\in A$, szereg
$e^a=\sum_{n\ge0}a^n/n!$ zbiega w normie i odpowiada
holomorficznemu rachunkowi funkcyjnemu funkcji $e^z$:
szereg Taylora zbiega jednostajnie na konturze, a
oszacowanie normy całki przez długość konturu i supremum
normy funkcji podcałkowej pozwala przejść do granicy.
Twierdzenie~\ref{thm:holomorphic-calculus31} daje
$\sigma(e^a)=e^{\sigma(a)}$.
Jeśli $q$ jest wielomianem bez zer na $\sigma(a)$,
a $p$ wielomianem, to $(p/q)(a)=p(a)q(a)^{-1}$.
\end{przyklad}

\subsection{Algebry \texorpdfstring{$C^*$}{C*} i ich przemienny model}

\begin{definicja}[Algebra $C^*$]
\label{def:cstar31}
\emph{Algebrą $C^*$} nazywamy zespoloną algebrę Banacha $A$
z odwzorowaniem $a\mapsto a^*$ spełniającym
$(a+b)^*=a^*+b^*$, $(\lambda a)^*=\bar\lambda a^*$,
$(ab)^*=b^*a^*$, $(a^*)^*=a$ oraz
$\|a^*a\|=\|a\|^2$. W tym rozdziale algebry $C^*$
są unitalne. Element $a$ jest \emph{samosprzężony}, gdy $a=a^*$,
a \emph{normalny}, gdy $a^*a=aa^*$.
\end{definicja}

\begin{przyklad}[Dwie algebry $C^*$]
Dla niezerowej zespolonej przestrzeni Hilberta $H$
algebra $\mathcal L(H,H)$ z operacją operatora
sprzężonego jest algebrą $C^*$.
Równość $\|T^*\|=\|T\|$ wynika z oszacowania
$\|T^*y\|=\sup_{\|x\|\le1}|\langle Tx,y\rangle|
\le\|T\|\|y\|$ i zastosowania tego oszacowania do $T^*$.
Stąd $\|T^*T\|\le\|T\|^2$, zaś dla $\|x\|\le1$
mamy $\|Tx\|^2=|\langle T^*Tx,x\rangle|\le\|T^*T\|$,
co po wzięciu supremum daje odwrotną nierówność.
Dla niepustej przestrzeni zwartej Hausdorffa $K$
algebra $C(K,\C)$ z operacją $f^*(x)=\overline{f(x)}$
jest przemienną algebrą $C^*$, ponieważ
$\|f^*f\|_\infty=\sup_{x\in K}|f(x)|^2=\|f\|_\infty^2$.
\end{przyklad}

\begin{lemat}[Norma i widmo elementu normalnego]
\label{lem:cstar-normal31}
Niech $A$ będzie unitalną algebrą $C^*$, a $a\in A$ elementem
normalnym. Wtedy $r(a)=\|a\|$.
Jeżeli $a=a^*$, to $\sigma(a)\subseteq\R$.
\end{lemat}
\begin{proof}
Inwolucja zachowuje jednostkę: dla każdego $b=a^*$
mamy $1_A^* b=(a1_A)^*=b$ i $b1_A^*=(1_Aa)^*=b$,
więc $1_A^*=1_A$. Z własności $C^*$ wynika
$\|b^*\|=\|b\|$ (dla $b=0$ równość jest oczywista):
$\|b\|^2=\|b^*b\|\le\|b^*\|\|b\|$,
a analogicznie po zamianie $b$ na $b^*$.
Dla samosprzężonego $h$ mamy
$\|h^2\|=\|h^*h\|=\|h\|^2$.
Jeżeli $a$ jest normalny, to
$\|(a^2)^*a^2\|=\|(a^*a)^2\|=\|a^*a\|^2$;
przez własność $C^*$ otrzymujemy $\|a^2\|=\|a\|^2$.
Ponieważ $a$ komutuje z $a^*$, każda jego potęga jest
normalna: $(a^m)^*=(a^*)^m$ komutuje z $a^m$.
Indukcja daje $\|a^{2^n}\|=\|a\|^{2^n}$.
Wzór na promień spektralny implikuje $r(a)=\|a\|$.

Jeśli $h=h^*$, to dla $t\in\R$ i $u_t=e^{ith}$
ciągłość antyliniowej inwolucji pozwala zastosować ją
wyraz po wyrazie w szeregu eksponenty, dając
$u_t^*=e^{-ith}$. Z własności homomorfizmu rachunku
holomorficznego (Twierdzenie~\ref{thm:holomorphic-calculus31})
i tożsamości $e^{itz}e^{-itz}=1$ wynika
$u_t^*u_t=u_tu_t^*=1_A$. Własność $C^*$ daje
$\|u_t\|^2=\|1_A\|=1$, więc $\|u_t\|=1$.
Na mocy twierdzenia o odwzorowaniu widma
$e^{it\lambda}\in\sigma(u_t)$ dla $\lambda\in\sigma(h)$.
Zatem $|e^{it\lambda}|\le1$ dla wszystkich $t\in\R$,
a po uwzględnieniu również $-t$ otrzymujemy
$\operatorname{Im}\lambda=0$.
\end{proof}

\begin{lemat}[Zespolone twierdzenie Stone'a--Weierstrassa]
\label{lem:stone-weierstrass31}
Niech $K$ będzie niepustą zwartą przestrzenią Hausdorffa, a
$B\subseteq C(K,\C)$ podalgebrą zawierającą funkcje stałe,
rozdzielającą punkty $K$ i zamkniętą na sprzężenie zespolone.
Wtedy $B$ jest gęsta w $C(K,\C)$ w normie supremum.
\end{lemat}
\begin{proof}
Najpierw ustalmy, że każdą rzeczywistą funkcję ciągłą na odcinku
można jednostajnie aproksymować wielomianami rzeczywistymi.
Po zmianie zmiennej wystarczy odcinek $[0,1]$. Dla $h\in C([0,1],\R)$
rozważamy wielomiany Bernsteina
$B_nh(t)=\sum_{k=0}^nh(k/n)\binom nk t^k(1-t)^{n-k}$.
Ponieważ wagi są nieujemne, sumują się do $1$ na mocy wzoru dwumianowego, a ich
średnia dla zmiennej $k/n$ wynosi $t$ i wariancja
$t(1-t)/n\le1/(4n)$ (z tożsamości
$k\binom nk=n\binom{n-1}{k-1}$ oraz
$k(k-1)\binom nk=n(n-1)\binom{n-2}{k-2}$ dla $n\ge2$,
z pominięciem wyrazów $k<2$; dla $n=1$ drugi moment
silniowy wynosi zero).
Rozdzielamy sumę na $|k/n-t|\le\delta$ i pozostałe
indeksy. W pierwszej części różnica wartości $h$ jest
ograniczona przez moduł ciągłości, w drugiej przez
$2\|h\|_\infty$, a suma wag drugiej części nie przekracza
$t(1-t)/(n\delta^2)$ na mocy oszacowania drugiego momentu.
Otrzymujemy, przy oznaczeniu
$\omega_h(\delta)=\sup\{|h(s)-h(t)|:s,t\in[0,1],\ |s-t|\le\delta\}$,
oszacowanie
$\|B_nh-h\|_\infty\le\omega_h(\delta)
+\|h\|_\infty/(2n\delta^2)$ dla każdego $\delta>0$.
Wybierając najpierw małe $\delta$, a następnie duże $n$,
otrzymujemy zbieżność jednostajną.

Niech $B_\R=\{f\in B:f(K)\subseteq\R\}$.
Dzięki zamknięciu na sprzężenie, część rzeczywista i urojona
dowolnego elementu $B$ należą do $B_\R$.
Oznaczmy przez $D$ domknięcie $B_\R$ w $C(K,\R)$.
Domknięcie $D$ jest algebrą rzeczywistą z jednostką,
ponieważ dodawanie i mnożenie są ciągłe w normie supremum.
Dla $u\in D$ aproksymujemy funkcję $s\mapsto|s|$
wielomianami na $[-\|u\|_\infty,\|u\|_\infty]$
(przypadek $u=0$ jest natychmiastowy). Ich złożenia
z $u$ należą do $D$ i zbiegają jednostajnie do $|u|$.
Zatem $|u|\in D$.
Stąd $\max(u,v)=(u+v+|u-v|)/2$ i
$\min(u,v)=(u+v-|u-v|)/2$ należą do $D$.

Ustalmy $h\in C(K,\R)$ i $\varepsilon>0$.
Dla $x\ne y$ wybieramy element $B$ o różnych wartościach
w $x,y$. Jego część rzeczywista lub urojona daje
$v\in B_\R$ z $v(x)\ne v(y)$. Kładziemy
$u_{xy}=h(x)+(h(y)-h(x))(v-v(x))/(v(y)-v(x))$.
Dla $x=y$ wybieramy stałą $h(x)$.
W obu przypadkach $u_{xy}(x)=h(x)$ i $u_{xy}(y)=h(y)$.
Dla ustalonego $x$ z ciągłości wynika, że dla każdego $y$
istnieje otoczenie $V_y$ punktu $y$, na którym
$u_{xy}<h+\varepsilon$. Po wybraniu skończonego podpokrycia
odpowiednia funkcja $g_x=\min_j u_{xy_j}\in D$ spełnia
$g_x<h+\varepsilon$ na całym $K$ i $g_x(x)=h(x)$.
W pewnym otoczeniu $U_x$ punktu $x$ zachodzi także
$g_x>h-\varepsilon$. Wybierając skończenie wiele takich
otoczeń pokrywających $K$ i biorąc $g=\max_j g_{x_j}\in D$,
otrzymujemy $|g-h|<\varepsilon$ na $K$.
Zatem $D=C(K,\R)$. Każdą funkcję zespoloną można rozłożyć
na część rzeczywistą i urojoną, więc $B$ jest gęsta w $C(K,\C)$.
\end{proof}

\begin{twierdzenie}[Przemienne twierdzenie Gelfanda--Naimarka]
\label{thm:commutative-cstar31}
Niech $A$ będzie przemienną unitalną algebrą $C^*$.
Wówczas transformacja Gelfanda jest izometrycznym izomorfizmem
algebr z inwolucją
$A\simeq C(\Delta(A),\C)$,
przy czym $\widehat{a^*}=\overline{\widehat a}$.
\end{twierdzenie}
\begin{proof}
Dla elementu samosprzężonego $h$ mamy
$\chi(h)\in\sigma(h)\subseteq\R$ dla każdego charakteru $\chi$.
Każdy $a$ rozkłada się jako $a=h+ik$ z
$h=(a+a^*)/2$, $k=(a-a^*)/(2i)$ samosprzężonymi.
Zatem $\chi(a^*)=\overline{\chi(a)}$.
Ponieważ w przemiennej algebrze każdy element jest normalny,
lemat~\ref{lem:cstar-normal31} i twierdzenie~\ref{thm:gelfand-spectrum31}
dają $\|\widehat a\|_\infty=r(a)=\|a\|$.
Izometria jest różnowartościowa, a jej obraz jest domknięty
przez zupełność $A$. Obraz zawiera stałe, ponieważ
$\widehat{\lambda1_A}=\lambda$. Rozdziela punkty:
dla $\chi_1\ne\chi_2$ istnieje $a$ z
$\chi_1(a)\ne\chi_2(a)$. Jest podalgebrą zamkniętą
na sprzężenie dzięki wykazanej tożsamości dla $a^*$.
Z lematu~\ref{lem:stone-weierstrass31} wynika,
że jest gęsty w $C(\Delta(A))$, a więc równy całej tej przestrzeni.
Zatem odwzorowanie Gelfanda jest na $C(\Delta(A),\C)$.
\end{proof}

\begin{uwaga}[Granica zakresu]
Twierdzenie~\ref{thm:commutative-cstar31} dotyczy algebr $C^*$
\emph{przemiennych}. Pełne twierdzenie Gelfanda--Naimarka o
wiernej reprezentacji dowolnej algebry $C^*$ na przestrzeni Hilberta
wymaga m.in. konstrukcji GNS i będzie osobnym rozszerzeniem.
Nie utożsamiamy też rachunku holomorficznego z pełnym
ciągłym rachunkiem funkcyjnym dla operatorów normalnych;
ten drugi zostanie opracowany w \hyperref[sec:spectral-hilbert32]{rozdziale~32}.
\end{uwaga}


\subsection{Dwa przykłady pokazujące znaczenie struktury algebry}

\begin{przyklad}[Charaktery algebry splotowej $\ell^1(\N_0)$]
\label{ex:l1-characters31}
Niech $A=\ell^1(\N_0)$ z mnożeniem splotowym opisanym
w przykładzie~\ref{ex:banach-algebras31}. Oznaczmy przez $e_n$ ciąg,
którego $n$-ty wyraz jest równy $1$, a wszystkie pozostałe są zerowe.
Wówczas $e_n=e_1^{*n}$ i $e_0=1_A$.
Jeżeli $\chi\in\Delta(A)$, to $z:=\chi(e_1)$ spełnia
$|z|\le\|e_1\|_1=1$. Przez multiplikatywność i ciągłość charakteru
otrzymujemy dla dowolnego $a=(a_n)\in A$
\[
 \chi(a)=\sum_{n=0}^\infty a_nz^n.
\]
Odwrotnie, dla każdego $|z|\le1$ ten szereg określa charakter:
jest absolutnie zbieżny i zachowuje splot dzięki możliwości zamiany
kolejności sum bezwzględnie zbieżnych. Zatem przestrzeń charakterów
$\Delta(A)$ jest homeomorficzna z domkniętym dyskiem jednostkowym:
odwzorowanie $z\mapsto\chi_z$ jest ciągłe w topologii
słabej-* dzięki jednostajnej zbieżności szeregu
$\sum a_nz^n$ na domkniętym dysku dla każdego $a\in A$,
jest bijekcją, a odwrotność $\chi\mapsto\chi(e_1)$
jest ciągła z definicji tej topologii; a transformacja Gelfanda przyporządkowuje
ciągowi $a$ funkcję $\widehat a(z)=\sum_{n\ge0}a_nz^n$.
W szczególności $\sigma_A(e_1)=\overline{\mathbb D}$,
choć $e_1$ nie jest elementem odwracalnym.
\end{przyklad}

\begin{przyklad}[Przesunięcie jednostronne — widmo niebędące zbiorem wartości własnych]
\label{ex:unilateral-shift31}
Na zespolonej przestrzeni Hilberta $H=\ell^2(\N)$ niech
$S(x_1,x_2,\ldots)=(0,x_1,x_2,\ldots)$.
Jest to izometria o normie $1$, więc $\sigma(S)\subseteq\overline{\mathbb D}$.
Dla każdego $|\lambda|<1$ operator sprzężony $S^*$ ma niezerowy
wektor własny $v=(1,\overline\lambda,\overline\lambda^2,\ldots)$
z wartością własną $\overline\lambda$.
Istotnie, $\sum_{n\ge0}|\lambda|^{2n}<\infty$,
a $S^*$ jest przesunięciem w lewo, więc
$(\lambda I-S)^*v=0$. Dla każdego $x\in H$
iloczyn skalarny $\langle(\lambda I-S)x,v\rangle$
jest zerowy. Zatem obraz $\lambda I-S$ jest zawarty
we właściwej domkniętej podprzestrzeni $\{v\}^\perp$,
nie jest gęsty i operator nie jest surjektywny.
Stąd $\lambda\in\sigma(S)$.
Z domkniętości widma dostajemy $\sigma(S)=\overline{\mathbb D}$.
Jednak $S$ nie ma żadnych wartości własnych: równanie $Sx=\lambda x$
dla $\lambda\ne0$ wymusza kolejno $x_1=0,x_2=0,\ldots$,
a dla $\lambda=0$ daje $x=0$ z różnowartościowości $S$.
Pokazuje to, dlaczego w ogólnej teorii operatorów nie należy
utożsamiać widma ze zbiorem wartości własnych, co było możliwe
poza zerem w twierdzeniu Riesza--Schaudera dla operatorów zwartych.
\end{przyklad}
